% Paper 2: Bitcoin as Unique Neutral Settlement: A Seven-Property Elimination
% Compiled from BGT-PAPER-2.md via pandoc + template assembly
% Build: xelatex BGT-PAPER-2.tex

\documentclass[12pt]{article}

% --- Fonts and encoding ---
\usepackage{fontspec}
\setmainfont{Latin Modern Roman}
\usepackage{unicode-math}
\setmathfont{Latin Modern Math}

% --- Page layout ---
\usepackage[margin=1in]{geometry}
\usepackage{setspace}
\onehalfspacing

% --- Math ---
\usepackage{amsmath}

% --- Theorem environments ---
\usepackage{amsthm}
\newtheorem{theorem}{Theorem}[section]
\newtheorem{proposition}[theorem]{Proposition}
\newtheorem{lemma}[theorem]{Lemma}
\newtheorem{corollary}[theorem]{Corollary}
\theoremstyle{definition}
\newtheorem{definition}{Definition}[section]
\newtheorem{assumption}{Assumption}
\theoremstyle{remark}
\newtheorem*{remark}{Remark}

% --- Tables ---
\usepackage{booktabs}
\usepackage{longtable}
\usepackage{array}
\usepackage{tabularx}
\usepackage{calc}

% --- Page numbers ---
\usepackage{fancyhdr}
\pagestyle{fancy}
\fancyhf{}
\fancyfoot[C]{\thepage}
\renewcommand{\headrulewidth}{0pt}

% --- Misc ---
\usepackage{enumitem}
\usepackage{hyperref}
\hypersetup{
  colorlinks=true, linkcolor=black, citecolor=black, urlcolor=blue,
  pdftitle={Bitcoin as Unique Neutral Settlement: A Seven-Property Elimination},
  pdfauthor={Sean Hash},
  pdfkeywords={Bitcoin, neutral settlement, asset comparison, digital scarcity, reserve assets, capital allocation, property elimination, game theory, network effects}
}
\usepackage{microtype}

% --- Suppress section numbering (manual numbers in headings) ---
\setcounter{secnumdepth}{0}

\begin{document}

% ============================================================
% TITLE BLOCK
% ============================================================

\begin{center}
{\LARGE\bfseries Bitcoin as Unique Neutral Settlement:\\A Seven-Property Elimination}

\bigskip

{\large Sean Hash}\\
\url{bitcoingametheory.com}\\
\href{mailto:sean@bitcoingametheory.com}{sean@bitcoingametheory.com}

\bigskip

February 2026\\[4pt]
{\normalsize Working Paper}

\medskip

{\small
\textbf{JEL Codes:} G11, G12, G15, D82, E42, E51, L14\\[4pt]
\textbf{Keywords:} Bitcoin, neutral settlement, asset comparison, digital scarcity, reserve assets, capital allocation, property elimination, game theory, network effects
}
\end{center}

\bigskip

% ============================================================
% ABSTRACT
% ============================================================

\begin{abstract}
\noindent
Which asset can serve as neutral settlement in a multipolar world? The
companion paper (Hash, 2026a) establishes that the payoff advantage of
Exit over Stay is monotone increasing in adoption, and derives seven
necessary properties: protocol security, neutrality, permissionless
access, cheap finality, absolute scarcity, informational security, and
adaptive resilience. This paper applies those properties. We conduct a
systematic elimination across seven asset classes --- fiat currencies,
sovereign bonds, equities, real estate, gold, alternative blockchain
tokens, and Bitcoin --- showing that each traditional class violates at
least one required property. Bitcoin satisfies all seven, not because it
is perfect on any single dimension, but because it is the only asset
with an empty capture surface. We support the elimination with empirical
data: \$170 billion in ETF assets under management within two years
(institutional adoption cascades), \$1 billion settled for under \$500
in fees (settlement efficiency), and total hashrate recovery within six
months of China's comprehensive ban (network resilience). The analysis
also explains why Bitcoin's position derives from an unreplicable
historical sequence rather than technical superiority alone, making the
network effect argument structural rather than contingent. The analysis
does not require fiat collapse --- Bitcoin operates as settlement
finality layer in parallel with fiat payment rails.
\end{abstract}

\section{1. Introduction}\label{introduction}

The Exit Game (Hash, 2026a) proves that the payoff advantage of moving
capital from capturable settlement systems to neutral settlement is
monotonically increasing in adoption. But the framework is agnostic
about \emph{which} asset serves as the Exit destination. It derives
seven necessary properties (P1-P7) and shows that the Exit advantage is
monotone increasing --- but the identity of the neutral settlement asset
is an empirical question, not a theoretical one.

This paper answers that question. We take P1-P7 as given and ask: which
existing asset satisfies all seven simultaneously? The answer is reached
by elimination, not by advocacy. Each asset class is tested against each
property. Six classes fail. One survives.

The structure is deliberate. By separating the game-theoretic argument
(Paper 1) from the asset identification (this paper), we make the logic
auditable: a reader who rejects our claim about Bitcoin need only show
that some alternative satisfies P1-P7, without engaging the game theory.
A reader who accepts Bitcoin's properties need only verify the
elimination table, without reconstructing the dominance proof.


\section{2. The Property Framework}\label{the-property-framework}

We summarize the seven properties derived in Hash (2026a), Section 3.3.
Each property blocks a specific attack class that would violate neutral
settlement (Definition 1: immune to seizure, debasement, and political
capture).

\subsection{2.1 Formal Specifications}\label{formal-specifications}

\textbf{P1 (Protocol Security).} Mining constitutes a Nash equilibrium;
attack cost exceeds expected gain. Formally: \emph{C}\_attack +
\emph{P}\_collapse \textgreater{} Δ\_attack, and Σ \emph{δ}\^{}\emph{t}
· E{[}Π\_honest{]} \textgreater{} Δ\_attack − \emph{C}\_attack
(following Budish, 2018 and Biais et al., 2019).

\textbf{P2 (Neutrality).} No issuer, no foundation, no governance
mechanism capable of unilateral rule changes.
\textbar{}\emph{CS}\textbar{} = 0.

\textbf{P3 (Permissionless Access).} Any agent can initiate settlement
without third-party permission. Pr(censor access \textbar{} scale) ≈ 0.

\textbf{P4 (Cheap Finality).} Settlement of \$1 billion for fees below
\$500, with mathematical finality in under 60 minutes.

\textbf{P5 (Absolute Scarcity).} Fixed supply with zero supply
elasticity: d\emph{S}/d\emph{P} = 0.

\textbf{P6 (Informational Security).} Custody is mathematical (private
key knowledge). Cost of seizure scales superlinearly with targets.

\textbf{P7 (Adaptive Resilience).} Protocol upgrades via consensus
without introducing governance capture.

\subsection{2.2 Necessity and
Sufficiency}\label{necessity-and-sufficiency}

\noindent\textbf{Proposition 1 (Necessity).} P1-P7 are individually necessary.
Removing any single property enables the corresponding attack vector.
(Proved in Hash, 2026a.)

\noindent\textbf{Proposition 2 (Sufficiency).} P1-P7 are jointly sufficient.
Candidate additional properties either reduce to P1-P7 (e.g.,
``privacy'' reduces to P6) or introduce non-structural criteria that
would create capture surface, violating P2.


\section{3. Elimination}\label{elimination}

\subsection{3.1 Method}\label{method}

For each asset class, we identify which properties are violated and the
mechanism of violation. An asset fails the test if it violates
\emph{any} single property, since P1-P7 are individually necessary.

\subsection{3.2 Results}\label{results}

{%
\begin{longtable}[]{@{}
  >{\raggedright\arraybackslash}p{(\linewidth - 4\tabcolsep) * \real{0.3023}}
  >{\raggedright\arraybackslash}p{(\linewidth - 4\tabcolsep) * \real{0.4419}}
  >{\raggedright\arraybackslash}p{(\linewidth - 4\tabcolsep) * \real{0.2558}}@{}}
\toprule\noalign{}
\begin{minipage}[b]{\linewidth}\raggedright
Asset Class
\end{minipage} & \begin{minipage}[b]{\linewidth}\raggedright
Properties Violated
\end{minipage} & \begin{minipage}[b]{\linewidth}\raggedright
Mechanism
\end{minipage} \\
\midrule\noalign{}
\endhead
\bottomrule\noalign{}
\endlastfoot
C1: Fiat/CBDCs & P2, P5 & Issuers can debase (supply elasticity
\textgreater{} 0) and surveil/censor transactions \\
C2: Sovereign Bonds & P5 & Negative real yield under financial
repression; denominated in debaseable currency \\
C3: Equities & P2 & Subject to regulatory capture (margin requirements,
trading halts, dilution) \\
C4: Real Estate & P3, P6 & Illiquid, physically seizable, subject to
property tax and eminent domain \\
C5: Gold & P4, P6 & Settlement friction of 3-8\% (insurance, transport,
assay); physically seizable \\
C6: Alt-L1 tokens & P2 & Governance capture via foundation treasuries,
VC voting blocs, or proof-of-stake concentration \\
C7: Bitcoin & None identified & --- \\
\end{longtable}
}

\subsection{3.3 Detailed Analysis}\label{detailed-analysis}

\textbf{C1: Fiat Currencies and CBDCs.} Central banks set monetary
policy to serve domestic objectives (Assumption 1). The supply schedule
is discretionary: d\emph{S}/d\emph{P} ≠ 0 by design. This violates P5
directly. CBDCs add programmable censorship (violating P3) and state
surveillance (violating P6). No fiat currency has maintained purchasing
power over a 50-year horizon against hard assets.

\textbf{C2: Sovereign Bonds.} Bonds are claims on future fiat. When the
underlying currency is debased, bondholders bear the loss through
negative real yields. Financial repression --- holding interest rates
below inflation --- is the observed mechanism. This is not a market
failure; it is policy working as intended. Real yields on 10-year US
Treasuries were negative for 26 of 36 months from 2020-2023 (Federal
Reserve Bank of St.~Louis, FRED series DGS10 minus T10YIE).

\textbf{C3: Equities.} Productive assets generate real returns but are
subject to regulatory capture. Governments can impose capital controls,
windfall taxes, forced sales, trading halts, and beneficial ownership
reporting requirements. The capture surface is large:
\textbar{}\emph{CS}\_equities\textbar{} ≫ 0.

\textbf{C4: Real Estate.} Property is the paradigmatic \emph{seizable}
asset. Eminent domain, property tax, and zoning regulation constitute
permanent capture surfaces. Physical location is known (violating P6).
Settlement requires title search, escrow, and legal process (violating
P4). Cross-border real estate settlement is measured in weeks, not
minutes.

\textbf{C5: Gold.} This is the hardest case and deserves extended
treatment. Gold has served as neutral settlement for millennia. It fails
on two properties: P4 (settlement friction) and P6 (seizure). Physical
gold settlement incurs 3-8\% costs including insurance, transport,
assay, and storage (BullionVault, 2025). The 3,000-16,000x efficiency
gap versus Bitcoin settlement reflects the fundamental difference
between physical and informational assets. Executive Order 6102 (1933)
demonstrated that gold is seizable at national scale.

\emph{The verification game.} Gold settlement embeds a verification game
that has no Bitcoin analogue. Model the interaction between a seller
(\emph{S}) delivering an asset claimed to be gold and a buyer (\emph{B})
who must decide whether to verify:

{%
\begin{longtable}[]{@{}lll@{}}
\toprule\noalign{}
& \emph{S}: Authentic & \emph{S}: Forge \\
\midrule\noalign{}
\endhead
\bottomrule\noalign{}
\endlastfoot
\emph{B}: Trust & (\emph{V} − \emph{P}, \emph{P}) & (−\emph{V}, \emph{V}
+ \emph{P}) \\
\emph{B}: Verify & (\emph{V} − \emph{P} − \emph{C\_V}, \emph{P}) &
(−\emph{C\_V}, −\emph{C\_F}) \\
\end{longtable}
}

where \emph{V} is the asset value, \emph{P} is the price, \emph{C\_V} is
verification cost, and \emph{C\_F} is forgery cost. The equilibrium
depends on the relationship between \emph{C\_V} and \emph{C\_F}.

For gold, \emph{C\_V} is \textbf{monotone increasing in forgery
sophistication}: spray-painted lead is caught by visual inspection
(\textasciitilde\$0), but tungsten-core bars require destructive assay
(2-5\% of bar value) or ultrasound testing (\$200-500/bar). Documented
cases include tungsten-core bars in commercial LBMA inventories and
spray-painted lead bars in retail markets. The US gold reserve at Fort
Knox has not been independently audited since 1953 --- even
sovereign-scale verification is prohibitively costly.

\noindent\textbf{Lemma (Verification Cost Asymmetry).} Let
\emph{C\_V\^{}\{gold\}}(\emph{s}) be the cost of verifying gold against
a counterfeiter of sophistication \emph{s} ∈ {[}0, 1{]}. Let
\emph{C\_V\^{}\{btc\}} be the cost of verifying a Bitcoin UTXO. Then:

\begin{enumerate}
\def\labelenumi{(\roman{enumi})}
\item
  \emph{C\_V\^{}\{gold\}}'(\emph{s}) \textgreater{} 0 --- gold
  verification cost increases with forgery sophistication.
\item
  \emph{C\_V\^{}\{btc\}} = \emph{ε} ≈ 0, independent of attacker
  sophistication (under Assumption 3, computational hardness).
\item
  Gold verification is non-persistent: each transfer resets the
  verification game (the bar could have been swapped). Bitcoin
  verification is persistent: a confirmed UTXO remains valid until
  spent.
\end{enumerate}

\emph{Proof.} (i) follows from the physical structure of gold: detecting
surface-level fakes (visual) costs less than detecting deep fakes
(destructive assay). More sophisticated forgeries require more invasive
--- and more expensive --- testing. (ii) follows from Assumption 3:
verifying a digital signature is computationally trivial regardless of
the attacker's resources (the attacker cannot produce a valid signature
without the private key). (iii) follows from the difference between
physical and informational assets: a physical bar's composition can
change between inspections; a UTXO's validity is determined by the
blockchain state, which is globally consistent and tamper-evident. ∎

The implication for autonomous agents is categorical: an AI agent
verifying gold must either trust a human intermediary (introducing
counterparty risk at \emph{r} = 0) or deploy physical inspection
infrastructure (humanoid robots, automated assaying machines) whose
calibration itself depends on human trust chains. Bitcoin verification
requires only a full node --- software the agent controls entirely.

\emph{The gold-ETF objection.} Gold ETFs (GLD, IAU) settle digitally at
low cost, arguably satisfying P4. But gold ETFs are claims on custodial
gold, not gold itself --- they introduce counterparty risk (the
custodian) and regulatory capture (the issuer). An ETF can be frozen,
diluted, or restructured by its sponsor. If we accept ETF wrappers as
satisfying P4, then by the same logic Bitcoin ETFs satisfy P4 --- and
the underlying asset remains the object of analysis, not the wrapper.
Gold's P4 and P6 failures are properties of gold-as-settlement, not
gold-as-financial-product.

\textbf{C6: Alternative L1 Tokens.} Every alternative blockchain
protocol has at least one of: a foundation treasury that constitutes a
governance capture surface, a venture capital allocation that creates
concentrated voting power, or a proof-of-stake mechanism in which wealth
concentration maps directly to protocol control. Ethereum's transition
to proof-of-stake in September 2022 made this explicit: the largest
stakers have proportional influence over block production and
transaction ordering.

\emph{The stablecoin objection.} Stablecoins (USDC, USDT) and tokenized
treasuries settle digitally at low cost and denominate in familiar
units, arguably satisfying P4. But stablecoins fail P2 and P5 by
construction. The issuer maintains a freeze function --- Tether has
frozen hundreds of addresses at the direction of law enforcement
agencies, totaling over \$1 billion as of 2025. This means
\textbar{}\emph{CS}\_stablecoin\textbar{} \textgreater{} 0: a single
entity can unilaterally alter settlement outcomes for any address.
Stablecoins also fail P5: they are denominated in fiat currency, so the
underlying asset debases when the reference currency debases --- the
holder bears debasement risk with none of the yield. The capture is not
hypothetical: proposed US stablecoin legislation (2025-2026) would cap
yield payments to holders below market rates, a provision secured
through banking industry lobbying to protect deposit bases. This is the
P2 failure mode in its clearest form: a political coalition altering the
asset's rules to serve incumbent interests, precisely the capture that
Definition 1 requires immunity from.

\textbf{C7: Bitcoin.} Bitcoin satisfies all seven properties. P1 is
supported by the mining Nash equilibrium (Biais et al., 2019). P2 holds
because there is no issuer, no foundation, and
\textbar{}\emph{CS}\textbar{} = 0. P3 is demonstrated by the network's
continued operation across every jurisdiction that has attempted to ban
it. P4 is empirically verified (\$1B settled for \textless\$500). P5 is
algorithmic (21 million cap, enforced by node consensus). P6 follows
from cryptographic custody. P7 is demonstrated by the soft fork upgrade
path (SegWit, Taproot) without governance capture.

\subsection{3.4 Coexistence}\label{coexistence}

The elimination does not require fiat collapse. Bitcoin operates as
superior collateral and settlement finality in parallel with fiat
payment rails. The cascade pressure (Hash, 2026a, Proposition 2) applies
to the reserve settlement function. Fiat retains payment and
unit-of-account roles --- potentially indefinitely. This is not a
revolutionary claim; it is a structural one.


\section{4. Empirical Support}\label{empirical-support}

Theory without data is speculation. The following observations are
consistent with the framework.

\subsection{4.1 Institutional Adoption}\label{institutional-adoption}

Spot Bitcoin ETFs accumulated \$61.5 billion in net inflows since
January 2024, reaching approximately \$170 billion in total AUM.
BlackRock's IBIT crossed \$97 billion in 435 days --- faster than any
ETF in history (The Block, 2025). This is consistent with the cascade
dynamics predicted by the Exit Game: institutional actors crossing their
adoption thresholds $p_i^*$ in clusters.

\subsection{4.2 Sovereign Behavior}\label{sovereign-behavior}

El Salvador holds 7,529 BTC with daily purchases and geothermal mining
(Bitcoin Treasuries, 2025). Singapore issued 13 digital payment token
licenses in 2024 (MAS, 2024). Argentina received \$91.1 billion in
cryptocurrency value amid 143\% inflation (Chainalysis, 2024). These
observations are consistent with sovereign defection from the Stay
coalition (Hash, 2026a, Theorem 2).

\subsection{4.3 Settlement Efficiency}\label{settlement-efficiency}

Bitcoin settles \$1 billion in value for fees typically under \$500 with
mathematical finality in under 60 minutes (Mempool.space, 2024). Gold
settlement incurs 3-8\% friction costs (BullionVault, 2025). This
3,000-16,000x gap is the empirical basis for the P4 elimination of gold.

\subsection{4.4 Network Resilience}\label{network-resilience}

China's comprehensive mining ban in 2021 resulted in total hashrate
recovery to all-time highs within six months (CCAF, 2022). Bitcoin's
hashrate reached approximately 700 EH/s in 2025, with no PoW competitor
exceeding 1\% of this security budget (CCAF, 2025).

\subsection{4.5 Historical Uniqueness}\label{historical-uniqueness}

Bitcoin's position derives from an unreplicable historical sequence: (1)
pure origin (no ICO, no pre-mine, no venture capital, no foundation
treasury --- Nakamoto, 2008), (2) founderless development (creator
departed, estimated 1.1 million BTC unspent since 2009 --- Lerner,
2019), (3) survival of six categories of existential crisis without
protocol failure, and (4) PoW hashrate dominance following Ethereum's
migration to proof-of-stake in 2022.

No competitor can replicate this sequence. As Huberman, Leshno, and
Moallemi (2021) observe, Bitcoin is a ``monopoly without a monopolist''
--- and the monopoly derives from the \emph{absence} of a monopolist,
not despite it. A new protocol with a known founder, a foundation
treasury, and venture backing fails P2 by construction, regardless of
its technical merits.


\section{5. Discussion}\label{discussion}

\subsection{5.1 The Gold Question}\label{the-gold-question}

Gold is the strongest competitor and deserves extended treatment. Gold
has served as neutral settlement for millennia. Its failure on P4 and P6
is not a deficiency of gold but a consequence of its physicality. The
information revolution created a new possibility: settlement that is
mathematical rather than physical, where custody requires knowledge
rather than location.

Krugman (2018) argues that Bitcoin lacks fundamental value because it
has no tether to economic activity. This is reasonable if you weight
intrinsic yield over settlement properties. The present analysis weights
structure --- and on structure, gold fails two properties that Bitcoin
satisfies.

Taleb (2021) makes a different case: that Bitcoin is fragile because it
depends on continuous mining, electricity, and internet infrastructure.
This is the steelman for P7 failure. Our response: seventeen years of
operation through six categories of existential crisis (exchange
failures, regulatory bans, mining exodus, protocol disputes, market
crashes, and sustained negative press) constitute empirical evidence for
P7, though not proof. The fee-transition risk (Carlsten et al., 2016)
remains the strongest version of Taleb's concern.

\subsection{5.2 Attack Survival}\label{attack-survival}

The elimination holds only if Bitcoin actually satisfies P1-P7. The most
serious challenge is sovereign attacks: can state actors ban, seize, or
suppress Bitcoin?

The steelman: a coordinated coalition --- G7 plus China ---
simultaneously criminalizing possession, enforcing through ISP-level
filtering, and imposing secondary sanctions. The defense rests on three
structural features. Informational security (P6) transforms seizure into
a key-management problem. Permissionless access (P3) redirects activity
rather than eliminating it --- as China's 2021 ban empirically
demonstrated. And the coalition itself is unstable: the first defector
captures fleeing capital, and coordination problems are precisely what
this paper series is about.

We regard F1 --- permanent single global coordinator --- as the most
plausible falsification path. Not because it is likely, but because it
is the only attack that does not face a game-theoretic self-undermining
dynamic.

\subsection{5.3 The Volatility
Objection}\label{the-volatility-objection}

Bitcoin's annualized volatility has exceeded 70\% in multiple calendar
years, with drawdowns exceeding 50\% from all-time highs occurring in
2014, 2018, 2022, and again in early 2026 (approximately 52\% decline
from cycle highs). This is the most common objection from traditional
finance practitioners and deserves direct treatment.

Volatility is not among P1-P7. This is deliberate: volatility is a
market phenomenon, not a structural property of the protocol. The Exit
Game model (Hash, 2026a, condition M2) posits that
\emph{σ\_B}'(\emph{p}) \textless{} 0 --- deeper markets reduce
volatility. The long-term trend is consistent with this: Bitcoin's
90-day realized volatility has declined from \textasciitilde150\% in
2011 to \textasciitilde45\% in 2024. However, the empirical evidence is
mixed. Sharp drawdowns persist at current adoption levels. The power law
channel (price following a power law of time since genesis) describes
the central tendency but does not constrain short-term deviations.

We therefore distinguish two claims. The \emph{structural} claim:
volatility affects adoption timing (through the \emph{λ\_i} ·
\emph{σ\_B}(\emph{p}) term) but not equilibrium structure. More
risk-averse actors have higher thresholds $p_i^*$ and wait longer. This
holds regardless of \emph{σ\_B}'s sign. The \emph{empirical} claim:
volatility will decrease with adoption as markets deepen. This is
plausible but unproven at current adoption levels, and the 2026 drawdown
is a direct counterexample to the naive version. M2 should be understood
as a long-run tendency, not a monotone empirical fact at every time
horizon.

\subsection{5.4 Settlement
vs.~Acceptance}\label{settlement-vs.-acceptance}

The elimination framework rests on a distinction that deserves explicit
treatment: \emph{settlement} is not \emph{acceptance}.

Settlement in Bitcoin is a cryptographic fact. A valid signature
transfers value from address A to address B. The protocol does not ask
who owns the addresses, whether the transaction serves a lawful purpose,
or whether the receiving party will be able to spend the output. At the
protocol layer, P3 (permissionless access) holds unconditionally.

Acceptance is a social fact. If address B is associated with a
sanctioned entity --- a designation maintained by OFAC, the EU, or other
authorities --- then exchanges and custodians will refuse to credit
incoming funds from B or any address that received value from B. The
coins are settled but economically impaired: their acceptance by
regulated counterparties depends on off-chain arbitration that the
protocol cannot control.

This creates a race condition. Mixing services, coinjoins, and
cross-chain bridges attempt to break the provenance chain. Compliance
tools (Chainalysis, Elliptic) attempt to trace it. The effectiveness of
both evolves continuously. Partial laundering creates a spectrum of
``cleanliness'' rather than a binary clean/tainted classification, and
gray-market exchanges accept coins that regulated exchanges reject.

For our analysis, the implication is: P3 guarantees settlement at the
protocol layer, but the \emph{economic utility} of that settlement
depends on the post-settlement acceptance environment. The elimination
holds for settlement functionality --- Bitcoin is the unique asset where
the \emph{protocol} imposes no barriers. Whether the \emph{ecosystem}
around the protocol adds barriers is a compliance question, not a
settlement one. This distinction is sharpened in Hash (2026c), where
autonomous agents face the compliance problem in its purest form.

\textbf{The Acceptance Game.} The settlement-acceptance distinction
produces a natural enforcement equilibrium without requiring
protocol-level censorship. Consider an actor receiving coins of
uncertain provenance. The expected payoff of blind acceptance is:

E{[}Accept{]} = \emph{V} − \emph{p} · \emph{d}(\emph{t}) · \emph{c}

where \emph{V} is the value of the coins, \emph{p} is the probability of
taint, \emph{d}(\emph{t}) is the probability of detection at time
\emph{t}, and \emph{c} is the cost of punishment (legal penalties,
reputational damage, asset seizure). Acceptance is rational when
\emph{V} \textgreater{} \emph{p} · \emph{d}(\emph{t}) · \emph{c}.

Three features distinguish this game from the cash analogue. First,
\emph{d}(\emph{t}) is monotonically increasing: Bitcoin's permanent
ledger means chain analysis capabilities improve over time and apply
retroactively to past transactions. Coins accepted today may be flagged
in 2030. Cash has no such property --- once spent, provenance is lost.
Second, regulated entities (exchanges, custodians) face \emph{c} values
that dominate \emph{V} for any realistic taint probability --- license
revocation, criminal liability, institutional collapse --- and these
entities control the fiat on/off ramps. Third, the resulting
``provenance discount'' on tainted coins reduces the economic payoff of
the underlying crime: you can steal Bitcoin, but you cannot easily
convert it at par value.

The equilibrium stratifies by actor type:

{%
\begin{longtable}[]{@{}
  >{\raggedright\arraybackslash}p{(\linewidth - 6\tabcolsep) * \real{0.1296}}
  >{\raggedright\arraybackslash}p{(\linewidth - 6\tabcolsep) * \real{0.3148}}
  >{\raggedright\arraybackslash}p{(\linewidth - 6\tabcolsep) * \real{0.3519}}
  >{\raggedright\arraybackslash}p{(\linewidth - 6\tabcolsep) * \real{0.2037}}@{}}
\toprule\noalign{}
\begin{minipage}[b]{\linewidth}\raggedright
Actor
\end{minipage} & \begin{minipage}[b]{\linewidth}\raggedright
\emph{d} (detection)
\end{minipage} & \begin{minipage}[b]{\linewidth}\raggedright
\emph{c} (punishment)
\end{minipage} & \begin{minipage}[b]{\linewidth}\raggedright
Filtering
\end{minipage} \\
\midrule\noalign{}
\endhead
\bottomrule\noalign{}
\endlastfoot
Regulated exchange & High & Very high & Mandatory \\
Individual (KYC'd) & Medium & High & Incentivized \\
AI agent & Medium & Low (no personhood) & Rational (near-zero analysis
cost) \\
Gray market & Low & Low & Minimal \\
\end{longtable}
}

Notably, the actors most capable of filtering (regulated exchanges) are
also the actors most severely punished for failure, creating robust
enforcement at the fiat conversion bottleneck. For autonomous agents,
filtering converges to near-perfection not from moral obligation but
from the trivial cost of on-chain provenance analysis relative to any
nonzero expected punishment.

The policy implication: protocol-level censorship (violating P2 and P3)
is unnecessary for enforcement, because the application layer produces
its own enforcement equilibrium through rational acceptance decisions.
Bitcoin's neutrality at the settlement layer is compatible with ---
indeed, strengthened by --- discretionary filtering at the acceptance
layer.

\subsection{5.5 Limitations}\label{limitations}

The elimination is as strong as the property framework. If P1-P7 are
incomplete --- if there exists an eighth necessary property that Bitcoin
violates --- the elimination fails. The sufficiency argument
(Proposition 2) addresses this, but sufficiency proofs are inherently
harder than necessity proofs because they require exhaustive enumeration
of alternatives.

\emph{On minimality and completeness.} We claim that P1-P7 are
individually necessary and jointly sufficient. We do not claim they are
\emph{minimal} --- that is, that no proper subset of P1-P7 suffices.
Formal minimality would require showing that for each property P\_k,
there exists an asset satisfying all P\_j (j ≠ k) that fails as neutral
settlement. The counterexamples are suggestive (Ethereum for P1, Solana
for P2, CBDCs for P3, gold for P4, commodity bonds for P5, Swiss banking
for P6, Litecoin for P7) but a rigorous minimality proof --- in the
tradition of axiomatic characterizations in social choice (Arrow, 1951)
and mechanism design (Gibbard, 1973) --- would require 7-10 additional
pages and is not attempted here. The relevant question for the
elimination is necessity + sufficiency: can an asset fail any single
property and still serve as neutral settlement? The answer is no, and
that is what the framework proves.

Seventeen years of data is brief for a monetary claim. Gold has
millennia. The response --- that the relevant properties are structural,
not historical --- is logically sound but empirically untested at the
timescales that matter. Carlsten et al.'s (2016) analysis of the
fee-only security budget is particularly relevant: the transition from
block subsidies to fee-only security is the largest untested structural
assumption, and it will not be resolved by theory alone.


\section{6. Conclusion}\label{conclusion}

The elimination is complete. Among seven asset classes tested against
seven necessary properties, six classes fail at least one property.
Bitcoin fails none. This is not a claim about superiority on any single
dimension --- gold is more historically tested, equities generate higher
short-term returns, fiat currencies are more liquid for daily payments.
The claim is narrower: Bitcoin is the only existing asset with an empty
capture surface that simultaneously satisfies all requirements for
neutral settlement.

The companion paper (Hash, 2026a) proves that the payoff advantage of
Exit over Stay is monotonically increasing. This paper identifies the
destination. A third paper (Hash, 2026c) extends the analysis to the
case where the convergence pressure is strongest: autonomous agents with
zero legal recourse.


\section{References}\label{references}

Arrow, K. J. (1951). \emph{Social Choice and Individual Values}. John
Wiley \& Sons.

Biais, B., Bisiere, C., Bouvard, M., and Casamatta, C. (2019). The
blockchain folk theorem. \emph{Review of Financial Studies}, 32(5),
1662-1715.

Bitcoin Treasuries (2025). Sovereign Bitcoin holdings tracker.
bitcointreasuries.net.

Budish, E. (2018). The economic limits of Bitcoin and the blockchain.
\emph{NBER Working Paper} No.~24717.

BullionVault (2025). Gold settlement cost analysis.

Cambridge Centre for Alternative Finance (2022). Cambridge Bitcoin
Electricity Consumption Index: China mining ban impact.

Cambridge Centre for Alternative Finance (2025). Cambridge Bitcoin
Electricity Consumption Index: Global hashrate data.

Carlsten, M., Kalodner, H., Weinberg, S. M., and Narayanan, A. (2016).
On the instability of Bitcoin without the block reward.
\emph{Proceedings of the 2016 ACM SIGSAC}, 154-167.

Chainalysis (2024). The 2024 Geography of Cryptocurrency Report.

Federal Reserve Bank of St.~Louis (2024). FRED Economic Data: 10-Year
Treasury minus inflation expectations (DGS10, T10YIE series).

Gibbard, A. (1973). Manipulation of voting schemes: A general result.
\emph{Econometrica}, 41(4), 587--601.

Huberman, G., Leshno, J. D., and Moallemi, C. (2021). Monopoly without a
monopolist. \emph{Review of Financial Studies}, 34(7), 3266-3302.

Krugman, P. (2018). Transaction costs and tethers: Why I'm a crypto
skeptic. \emph{The New York Times}, July 31.

Lerner, S. (2019). The return of the deniers and the revenge of Patoshi.

Mempool.space (2024). Bitcoin transaction fee and settlement data.

Monetary Authority of Singapore (2024). Digital Payment Token Licensing
Statistics.

Nakamoto, S. (2008). Bitcoin: A peer-to-peer electronic cash system.

Hash (2026a). Bitcoin exit dominance in monetary coordination games.
Working paper, bitcoingametheory.com.

Hash (2026c). Settlement at zero trust: Bitcoin and autonomous economic
agents. Working paper, bitcoingametheory.com.

Taleb, N. N. (2021). Bitcoin, currencies, and fragility.
\emph{Quantitative Finance}, 21(8), 1249-1255.

The Block (2025). Bitcoin ETF AUM tracker.


\end{document}
